\clearpage
\textbf{The retained pre-run-9 64 KiB points are reproducible.} The retained c16 cases reproduce the high Native latency, not a particular cause. Table~\ref{tab:spike} shows each population separately. Native per-put barrier p50 spans \N{B:spike-summary:B_minimum:value|,.1f}--\N{B:spike-summary:B_maximum:value|,.1f}~\us; Native page-cache p50 spans \N{B:spike-summary:P_minimum:value|.1f}--\N{B:spike-summary:P_maximum:value|.1f}~\us. These are observed ranges, not confidence intervals or sustained-capacity results. Run-9 grouped-barrier controls are separate, not additions to these historical ranges.

\begin{table}[H]\centering\fitinput{tab-spike.tex}
\caption{Individual 64 KiB c16 successful-put API p50 values, recomputed from raw samples; \N{B:spike-cases:M_0379-native-65536-write-concurrent-primary_rf1-r0:n|,.0f} puts per population. B: Native per-put all-file/directory barrier, one process/SSD. P: Native page-cache acknowledgement, no per-put foreground fsync or per-put persistence guarantee. S: owner verbs synchronous covering persistence, five owners. V: Valkey RDMA AOF-always. Cache state is not controlled; coordinator cache is off in the distributed KV cases. N3: 3 November campaign; M: full matrix; R7/R8: focused fleets. R8 K/N/W are different profiles, not extra independent repetitions; columns retain each profile's repetition index. O/D: 4 November manual read/delete exports; Q: queue campaign. Copied baselines in manual exports are not proved independent fleets. Full case ids, sample/archive paths and contracts are in \texttt{spike-cases.csv}. No populations are pooled.}\label{tab:spike}
\end{table}

The full-matrix Native barrier distribution has p10/p50/p99 \N{B:spike-cases:M_0379-native-65536-write-concurrent-primary_rf1-r0:p10_us|,.1f}/\N{B:spike-cases:M_0379-native-65536-write-concurrent-primary_rf1-r0:p50_us|,.1f}/\N{B:spike-cases:M_0379-native-65536-write-concurrent-primary_rf1-r0:p99_us|,.1f}~\us. Only \N{B:spike-cases:M_0379-native-65536-write-concurrent-primary_rf1-r0:above_twice_pct|.2f} percent exceed twice its median. Its equal-count, dispatch-ordered quarter medians are \N{B:spike-cases:M_0379-native-65536-write-concurrent-primary_rf1-r0:q1_p50_us|,.1f}, \N{B:spike-cases:M_0379-native-65536-write-concurrent-primary_rf1-r0:q2_p50_us|,.1f}, \N{B:spike-cases:M_0379-native-65536-write-concurrent-primary_rf1-r0:q3_p50_us|,.1f} and \N{B:spike-cases:M_0379-native-65536-write-concurrent-primary_rf1-r0:q4_p50_us|,.1f}~\us. The main body is slow, not only a few outliers; the quarters show no monotonic warm-up decline.

The full-matrix page-cache case also has a slow main body: p10/p50/p99 \N{B:spike-cases:M_0490-native-65536-write-concurrent-native_async_write-r0:p10_us|.1f}/\N{B:spike-cases:M_0490-native-65536-write-concurrent-native_async_write-r0:p50_us|.1f}/\N{B:spike-cases:M_0490-native-65536-write-concurrent-native_async_write-r0:p99_us|.1f}~\us. Its last quarter is slower than its first; this does not identify a cause.

The following one-barrier-per-put policy describes the retained Native per-put arm, not the separate run-9 grouped Native barrier.
The barrier policy is known: one library barrier after every put, without an application-level covering group. \inference\ The serialized-service model is consistent with the earlier c1/c16 ratios, but kernel sync work can overlap. The page-cache jump from 8 to 64 KiB remains unexplained: \N{J:pagecache_jump} times the p50 for eight times the payload. Worker identity and per-put WAL offsets are absent from \path{samples.csv}; fragment-boundary clustering cannot be tested. The API clock includes engine frame encoding, CRC and WAL copying, but excludes harness value generation and oracle bookkeeping. Sources: \path{runner.rs:2655--2808}, \path{backends/native.rs:173--179,295--300}, \path{db.rs:579--594,2157--2167} and \path{crc.rs:17--29}; paths are relative to the benchmark or engine source directories.

\clearpage
\textbf{Local structure test (laptop, not the cloud).} This later audit used IEEE CRC32, which matches the engine's checksum polynomial, not its full implementation. The machine ran kernel \texttt{7.2.6}, ext4 on local NVMe, not cloud Local SSD. The local c16 slowdown appears at 8 KiB as well as 64 KiB; it does not reproduce the cloud's size-specific page-cache jump. \inference\ Write-prefault and \texttt{pwrite} still slow down under concurrency, so neither page faults nor copy contention is rejected. The grouped ticket model was not uniformly faster: its 64 KiB c16 p50 was \N{B:local-structure-test:grouped_data_sync_64_c16_r1:p50_ms|.2f}/\N{B:local-structure-test:grouped_data_sync_64_c16_r2:p50_ms|.2f} ms in runs 1/2, against \N{B:local-structure-test:data_sync_64_c16_r1:p50_ms|.2f} ms for per-put data sync. It is not evidence for a grouped-barrier speedup. The cloud step's cause stays open.

\begin{table}[H]\centering\fitinput{tab-local-structure.tex}
\caption{Local structure test: transcribed from the audit report; raw samples not retained. Each cell lists IEEE run 1; run 2 separately; -- means not measured. Fresh mapping, \texttt{MAP\_POPULATE}, write-prefault, \texttt{pwrite} and anonymous memory have no per-put foreground persistence guarantee. Write-prefault writes every page and data-syncs before timing; \texttt{MAP\_POPULATE} is not that control. Data sync means per-put \texttt{fdatasync}; grouped sync is a covered-ticket model. Neither is the engine's complete file-and-directory barrier. Finite-case rates are not capacity estimates; per-call occupancy cannot be checked without retained samples. Source and run labels are in \texttt{local-structure-test.csv}.}\label{tab:local-structure}
\end{table}

\begin{minipage}{\linewidth}
\subsubsection{Fair Native comparison: compact subset measured}\label{sec:fair-native}
Run 9 measures the one-Db Native and one-owner TCP/verbs subset at equal total writers \N{B:fair-native-plan:G14N:writer0|.0f}, \N{B:fair-native-plan:G14N:writer1|.0f} and \N{B:fair-native-plan:G14N:writer2|.0f}. Native page-cache, per-put barrier and grouped barrier remain distinct from owner covering persistence (Section~\ref{sec:rwc}). The five-Db Native and five-owner rows below remain unmeasured design inventory; page-cache acknowledgement is not durable.

\begin{table}[H]\centering\fitinput{tab-fair-native.tex}
\caption{Frozen design inventory, not counts of completed run-9 cases. Run 9 includes G14N/G14O1 and three separate G14P diagnostics; G14S and G14O5 were not run. Instances, SSD location/count, network/security, sync scope and routing remain unmatched.}\label{tab:fair-native}
\end{table}

\texttt{fair\_native} has \N{B:fair-native-plan:fair_native:cases|.0f} designed cases; \texttt{parity} has \N{B:fair-native-plan:parity:cases|.0f}, including \N{B:fair-native-plan:G13W:cases|.0f} G13W floors. Run 9 executed the authorized \texttt{remote\_rw\_compact} subset, not either entire larger selection. Frozen digests in \texttt{fair-native-plan.csv} remain unchanged. A five-SSD Native arm remains unimplemented; further launches need separate approval.
\end{minipage}

